% ================================================================
%  第二章  行业背景和市场分析
%  写作建议：突出市场总量（千亿规模）、问题难度（国际/国内领先）
%  和前沿方向（AI、区块链等）。分析与所做的事紧密相关，多用数据/案例。
% ================================================================

\chapter{行业背景和市场分析}

% ----------------------------------------------------------------
% 2.1 该行业发展程度与现状动态
%   · 行业发展阶段定位（萌芽期/成长期/成熟期/整合期）
%   · 近2~3年行业内重大事件、技术突破、商业里程碑
%   · 国家/地方政策对本行业的最新导向与支持力度
%   · 行业主要参与者格局（头部企业、新兴玩家、投融资动向）
% ----------------------------------------------------------------
\section{该行业发展程度与现状动态}
% TODO:
%   - 行业阶段判断：目前处于XXX阶段，主要特征为……
%   - 近期重大事件：如「XXX企业于202X年完成XXX亿融资」「XXX政策于202X年落地」
%   - 政策环境：国家/省级政策名称及核心内容，对本行业的利好方向
%   - 竞争格局：行业CR5集中度，头部玩家规模，是否存在"强者恒强"趋势

% ----------------------------------------------------------------
% 2.2 行业总销售额与发展趋势
%   · 提供可信来源的行业规模历史数据（近3~5年）
%   · 引用预测机构对未来3~5年的增长预期（CAGR）
%   · 对增长驱动力进行拆因分析
% ----------------------------------------------------------------
\section{行业总销售额与发展趋势}
% TODO:
%   - 历史数据：202X年行业规模XXX亿元（来源：XXX机构）
%             202X年达到XXX亿元，同比增长XX%
%   - 趋势预测：预计2028年市场规模将达XXX亿元，CAGR约XX%（来源：XXX）
%   - 增长驱动：
%       · 需求驱动（人口结构变化/消费升级/数字化转型）
%       · 供给驱动（技术成本下降/新产品品类爆发）
%       · 政策驱动（专项资金/税收优惠/强制性标准）
%   % 建议配合折线图或柱状图，图表放在 image/chap02/ 目录下

% ----------------------------------------------------------------
% 2.3 经济与政府对该行业的影响
%   · 宏观经济走势如何影响本行业（顺周期/逆周期/独立增长）
%   · 政府补贴、采购、监管政策的具体影响
%   · 国际环境变化（汇率、供应链、出口限制等）的潜在影响
% ----------------------------------------------------------------
\section{经济与政府对该行业的影响}
% TODO:
%   - 经济影响：GDP增速/消费意愿对本行业市场需求的传导机制
%   - 政府支持：列举具体政策文件名称、补贴标准、扶持对象
%     （如："工业和信息化部《XXX行动计划（2024-2026年）》"）
%   - 监管风险：本行业面临的合规要求（资质/许可证/数据合规等）

% ----------------------------------------------------------------
% 2.4 市场细分
%   · 按多个维度对市场进行切分（用户特征/地理/使用场景/价格带）
%   · 给出每个细分市场的规模估算和增长性评价
%   · 识别哪些细分市场尚未被充分服务（机会所在）
% ----------------------------------------------------------------
\section{市场细分}
% TODO:
%   按【用户类型】细分：
%     - B端（企业客户）：规模XXX亿，主要需求为XXX
%     - C端（个人消费者）：规模XXX亿，主要需求为XXX
%     - G端（政府机构）：规模XXX亿，主要需求为XXX
%   按【地理区域】细分：一线城市/新一线城市/下沉市场/境外市场
%   按【使用场景】细分：细分场景A（规模/增长率）、场景B……
%   % 可使用饼图或矩阵图展示各细分市场占比

% ----------------------------------------------------------------
% 2.5 目标市场
%   · 在上述细分市场中，明确选择进入哪一个（或哪几个）
%   · 说明选择此目标市场的理由（需求强烈、竞争空白、资源匹配）
%   · 量化目标市场的规模（SAM）及本项目可获取的份额（SOM）
% ----------------------------------------------------------------
\section{目标市场}
% TODO:
%   - 目标市场定义：本项目聚焦于"XXX细分市场"，目标客户画像为：
%       年龄/职业/地域/行为特征/消费能力/痛点描述
%   - 目标市场规模（SAM）：约XXX亿元（计算方式：XXX×单价XXX×转化率XX%）
%   - 短期可获取份额（SOM）：预计第1年占据XX%，第3年达到XX%

% ----------------------------------------------------------------
% 2.6 市场定位
%   · 在竞争版图上确定本项目的坐标（价格×性能/高端×低端等二维图）
%   · 用 1~2 句话明确价值主张（Value Proposition）
%   · 与现有竞品的差异化机会分析，说明"空位"在哪里
% ----------------------------------------------------------------
\section{市场定位}
% TODO:
%   - 定位陈述："针对 [目标客户]，本项目提供 [核心产品/服务]，
%     与 [主要竞品] 相比，我们的独特优势在于 [差异化优势]。"
%   - 定位坐标图（可插入双轴散点图，横轴=价格，纵轴=技术/体验）
%     坐标图图片命名建议：image/chap02/positioning-map.png
%   - 竞争空白说明：现有玩家集中于XXX象限，本项目占据XXX空缺
\includegraphics[width=0.6\textwidth]{image/chap02/positioning-map}
% 注：如图片尚未准备好，请注释掉上面一行　
